\documentclass[11pt,a4paper]{article}
\usepackage[margin=2.5cm]{geometry}
\usepackage{times}
\usepackage{latexsym}
\usepackage{amsmath}
\usepackage{amssymb}
\usepackage{booktabs}
\usepackage{multirow}
\usepackage{graphicx}
\usepackage{xcolor}
\usepackage{natbib}
\usepackage{hyperref}
\usepackage[T1]{fontenc}
\usepackage[utf8]{inputenc}

\newcommand{\sys}{OntoCoDec}
\newcommand{\dr}{DR}
\newcommand{\rv}{RV}
\newcommand{\et}{ET}
\newcommand{\cf}{CF}

\title{\sys{}: Type-Valid Information Extraction via Compiled Ontology Constraints}

\author{Anonymous}

\begin{document}
\maketitle

\begin{abstract}
Large language models applied to information extraction produce systematic ontology constraint violations---generating triples where entity types are incompatible with relation domain/range requirements.
We present the first systematic study of these violations across two IE benchmarks (SciERC, DocRED) and two model families (Qwen3-8B, Llama-3.1-8B), finding that domain/range violations affect 13--26\% of extracted triples and concentrate on relation types with tight constraints.
Schema-informed prompting reduces violations by 32--62\% but is fundamentally unreliable: specific relation types exhibit \emph{increased} violation rates under prompt intervention.
We introduce \sys{}, a framework that automatically compiles IE ontology schemas into two composable constraint layers for autoregressive decoding: (1)~a grammar layer enforcing JSON structure and type vocabulary via context-free grammars, and (2)~a semantic layer enforcing cross-field domain/range constraints via a stateful entity registry and dynamic logit masking.
The grammar layer provides an \emph{enabling prerequisite} for the semantic layer: by ensuring all generated relation names belong to the valid vocabulary, it allows the semantic layer to look up and enforce domain/range rules.
On SciERC, \sys{} reduces domain/range violations by 94\% (15.4\%$\to$0.93\%) with F1-neutral impact ($p{=}0.58$).
On DocRED with Llama-3.1-8B, violations drop by 90.2\% (26.1\%$\to$2.55\%) with a significant F1 improvement ($p{<}0.001$).
All experiments use zero-shot prompting with zero code modification across models.
\end{abstract}

\section{Introduction}

LLM-based information extraction has shown strong zero-shot and few-shot capabilities across diverse domains \cite{Wei2023ZeroShotIE,Wadhwa2023Revisiting}.
However, when prompted to extract structured knowledge according to a predefined ontology, LLMs frequently produce \emph{ontology constraint violations}: triples where entity types are incompatible with the domain or range specifications of the predicted relation.
For instance, a model may extract the triple (\textit{CNN}, \textsc{Method}, \textsc{Evaluate-for}, \textit{image classification}, \textsc{Task}), violating the constraint that \textsc{Evaluate-for} requires a \textsc{Metric} in its range position.

These violations are not rare edge cases.
Our systematic measurement across two benchmarks reveals that 13--26\% of extracted triples contain domain/range violations, with rates reaching 42\% for relation types with tight constraints (\S\ref{sec:landscape}).
Prompt-level intervention---explicitly describing constraints in the system prompt---reduces violations overall but is \emph{fundamentally unreliable}: on SciERC, the \textsc{Evaluate-for} relation shows a 5.8 percentage point \emph{increase} in violations under schema-informed prompting.

We introduce \sys{} (\textbf{Onto}logy-\textbf{Co}nstrained \textbf{Dec}oding), a framework that compiles IE ontology schemas into decoding-time constraints.
\sys{} operates through two composable layers:
\begin{enumerate}
\item A \textbf{grammar layer} compiles the ontology's entity types and relation names into a context-free grammar (via XGrammar) that constrains output to valid JSON with legal type vocabularies.
\item A \textbf{semantic layer} maintains a stateful entity registry during generation and dynamically masks logits to enforce domain/range compatibility between entity types and relation arguments.
\end{enumerate}

The two layers exhibit an \emph{enabling dependency}: the grammar layer ensures every generated relation name belongs to the valid vocabulary, which is a prerequisite for the semantic layer to look up domain/range rules.
Without grammar enforcement, the semantic layer encounters out-of-vocabulary relation names and cannot enforce constraints---explaining why semantic-only enforcement achieves only $\sim$25\% violation reduction versus 94\% with the full system.

We evaluate \sys{} on SciERC \cite{Luan2018SciERC} (7 relations) and DocRED \cite{Yao2019DocRED} (96 relations) with two model families (Qwen3-8B, Llama-3.1-8B), all in zero-shot settings with zero code modification across models.
Our contributions are:
\begin{itemize}
\item The first systematic measurement of ontology constraint violations in LLM-based IE, revealing that violation rates are driven by constraint tightness rather than ontology size (\S\ref{sec:landscape}).
\item A composable two-layer constraint framework with principled enabling dependency, achieving 45--99\% domain/range violation reduction at zero or positive F1 cost (\S\ref{sec:method}, \S\ref{sec:experiments}).
\item A decidability taxonomy classifying ontology constraints by their amenability to forward-only autoregressive enforcement (\S\ref{sec:decidability}).
\end{itemize}

\section{Background}
\label{sec:background}

\paragraph{IE Ontology Structure.}
An IE ontology $\mathcal{O} = (\mathcal{E}, \mathcal{R}, \mathcal{C})$ defines entity types $\mathcal{E}$, relation types $\mathcal{R}$, and constraints $\mathcal{C}$.
For each relation $r \in \mathcal{R}$, the domain constraint $\text{dom}(r) \subseteq \mathcal{E}$ specifies valid head entity types and the range constraint $\text{rng}(r) \subseteq \mathcal{E}$ specifies valid tail entity types.
A triple $(e_h, t_h, r, e_t, t_t)$ is \emph{type-valid} iff $t_h \in \text{dom}(r)$ and $t_t \in \text{rng}(r)$.

\paragraph{Constraint Tightness.}
We define the tightness $\tau(r)$ of relation $r$ as $1 - \frac{|\text{dom}(r)| \cdot |\text{rng}(r)|}{|\mathcal{E}|^2}$, the fraction of entity type pairs excluded by the constraint.
Relations with $\tau \approx 1$ (e.g., \textsc{Compare} requiring identical types in both arguments) are tightly constrained; those with $\tau \approx 0$ accept most type combinations.

\paragraph{Grammar-Constrained Decoding.}
Grammar-constrained decoding restricts autoregressive generation to strings accepted by a formal grammar $G$, typically a context-free grammar.
At each decoding step, a token mask $M_G$ is computed such that $M_G[v] = 1$ iff token $v$ can extend the current prefix to a string in $L(G)$.
XGrammar \cite{Dong2024XGrammar} provides efficient CFG-constrained decoding for LLMs with precomputed adaptive token masks.

\section{Violation Landscape}
\label{sec:landscape}

Before developing constraint enforcement, we systematically measure ontology violations in LLM-based IE.
We evaluate Qwen3-8B in a zero-shot setting on SciERC (7 relation types, 6 entity types, 551 samples) and DocRED (96 relation types, 7 entity types, 998 documents).

\paragraph{Violation Taxonomy.}
We measure four violation types:
(1)~\textbf{Entity type violations} (\et{}): generated entity type $\notin \mathcal{E}$;
(2)~\textbf{Relation vocabulary violations} (\rv{}): generated relation name $\notin \mathcal{R}$;
(3)~\textbf{Domain/range violations} (\dr{}): entity types incompatible with relation constraints;
(4)~\textbf{Cross-field consistency violations} (\cf{}): same entity assigned different types across contexts.

\begin{table}[t]
\centering
\small
\begin{tabular}{lrrrr}
\toprule
& \multicolumn{2}{c}{\textbf{SciERC}} & \multicolumn{2}{c}{\textbf{DocRED}} \\
\cmidrule(lr){2-3} \cmidrule(lr){4-5}
\textbf{Violation} & \textbf{Rate} & \textbf{$n$} & \textbf{Rate} & \textbf{$n$} \\
\midrule
ET & 0.4\% & 2474 & 7.5\% & 19099 \\
RV & 1.4\% & 1724 & 28.9\% & 16540 \\
DR & 15.4\% & 1724 & 18.7\% & 16540 \\
CF & 0.08\% & 2474 & 0.08\% & 19099 \\
\bottomrule
\end{tabular}
\caption{Violation rates under minimal zero-shot prompting. DR rates are comparable across datasets despite 14$\times$ difference in ontology size (7 vs.\ 96 relations), supporting the tightness hypothesis.}
\label{tab:violation-overview}
\end{table}

\paragraph{Violations Are Driven by Tightness, Not Ontology Size.}
Table~\ref{tab:violation-overview} shows that SciERC (7 relations) and DocRED (96 relations) produce comparable DR rates (15.4\% vs.\ 18.7\%) despite a 14$\times$ difference in ontology complexity.
Within SciERC, per-relation DR ranges from 0.76\% (\textsc{Used-for}, $\tau{=}0.56$) to 42.4\% (\textsc{Compare}, $\tau{=}0.97$).
This 56$\times$ spread demonstrates that violation rate is determined by the constraint tightness of individual relations, not by the number of relations in the ontology.

\paragraph{Schema-Informed Prompting Is Unreliable.}
We test a schema-informed prompt that explicitly describes domain/range constraints for each relation.
On SciERC, this reduces overall DR from 15.4\% to 5.78\% ($-$62.4\%), but the reduction is \emph{unreliable per-relation}: \textsc{Feature-of} drops from 33.5\% to 11.2\% ($-$22.3pp), while \textsc{Evaluate-for} \emph{increases} from 24.1\% to 29.9\% ($+$5.8pp).
On DocRED, the schema-informed prompt achieves only 32\% DR reduction (vs.\ 62\% on SciERC), confirming that prompt-level intervention scales poorly with ontology complexity.

The \textsc{Evaluate-for} regression illustrates a \emph{prior-constraint conflict}: the schema prompt interferes with the model's learned type associations, destabilizing type assignment for relations where the model's distributional prior conflicts with ontology constraints.
This demonstrates that prompt-level intervention is not merely weak but \emph{mechanistically unreliable}---it can be actively counterproductive for specific constraint configurations.

\section{\sys{} Framework}
\label{sec:method}

\sys{} automatically compiles an IE ontology $\mathcal{O}$ into decoding-time constraints.
The compilation produces two constraint layers that compose via bitwise AND of token masks at each decoding step.

\subsection{Grammar Layer}

The grammar layer compiles $\mathcal{O}$ into a context-free grammar $G_\mathcal{O}$ that enforces:
(1)~valid JSON structure for IE output (entities list + relations list),
(2)~entity types restricted to $\mathcal{E}$,
(3)~relation names restricted to $\mathcal{R}$.

The grammar is compiled once per ontology and loaded into XGrammar for efficient mask computation.
This layer eliminates ET and RV violations by construction: every generated entity type and relation name is guaranteed to belong to the valid vocabulary.

\subsection{Semantic Layer}

The semantic layer enforces domain/range constraints through a \texttt{TypeConstraintProcessor} that operates during autoregressive generation:

\begin{enumerate}
\item \textbf{Entity Registry.} As entity mentions are generated (in the entities section of the output), the processor records each entity's name-to-type mapping in a registry $\mathcal{M}: \text{name} \to \text{type}$.
\item \textbf{Domain/Range Masking.} When generating the head or tail entity of a relation triple, the processor retrieves the current relation $r$ from the generation context, looks up valid types $\text{dom}(r)$ or $\text{rng}(r)$, and masks logits to allow only entity names whose registered types satisfy the constraint.
\end{enumerate}

\subsection{Enabling Dependency}
\label{sec:enabling}

The two layers are not independently effective---they exhibit an \emph{enabling dependency}.
The semantic layer requires a valid relation name to look up domain/range constraints in $\mathcal{O}$.
Without grammar enforcement, out-of-vocabulary relation names cause semantic lookup failures, and the semantic layer cannot enforce constraints.
Formally, for a generated relation name $\hat{r}$:
\[
\hat{r} \notin \mathcal{R} \implies \text{dom}(\hat{r}) = \text{rng}(\hat{r}) = \mathcal{E} \quad \text{(no constraint)}
\]
This explains why semantic-only enforcement achieves only $\sim$25\% DR reduction across both datasets: approximately 1--30\% of generated relation names are out-of-vocabulary (depending on model and dataset), and those triples escape semantic enforcement entirely.

When the semantic mask $M_{\text{sem}}$ and grammar mask $M_{\text{gram}}$ produce an empty intersection ($M_{\text{sem}} \cap M_{\text{gram}} = \emptyset$), the system falls back to grammar-only constraints, maintaining structural validity as a guaranteed floor.

\subsection{Decidability Taxonomy}
\label{sec:decidability}

We classify ontology constraints by their amenability to forward-only enforcement:
\begin{itemize}
\item[\textbf{(a)}] \textbf{Purely local}: entity type $\in \mathcal{E}$, relation name $\in \mathcal{R}$. Enforceable by grammar alone.
\item[\textbf{(b)}] \textbf{Prefix-decidable with state}: domain/range compatibility, given the entity registry accumulated during generation. Enforceable by the semantic layer.
\item[\textbf{(c)}] \textbf{Truly non-local}: forward coreference---an entity mentioned later may resolve the type of an earlier mention. Out of scope for single-pass forward-only decoding; responsible for residual violations.
\end{itemize}

\subsection{S1 and S2 Decoding Strategies}

We evaluate two strategies:
\textbf{S1 (single-pass)}: both layers active throughout generation. This is the primary recommended configuration.
\textbf{S2 (two-pass)}: Pass~1 extracts entities with grammar-only constraints; Pass~2 generates relations with an input-dependent grammar compiled from Pass~1 entities.
S2 is an architectural variant demonstrating framework composability; S1 is strictly preferred on all metrics (\S\ref{sec:ablation}).

\section{Experiments}
\label{sec:experiments}

\subsection{Setup}

We evaluate on SciERC \cite{Luan2018SciERC} (551 samples, 7 relation types, 6 entity types) and DocRED \cite{Yao2019DocRED} (998 documents, 96 relation types, 7 entity types) with Qwen3-8B and Llama-3.1-8B in zero-shot settings.
F1 is computed using the all\_set method: out-of-vocabulary predictions are included in the denominator, so eliminating OOV names via grammar enforcement directly reduces false positives.
Statistical significance is assessed via bootstrap resampling (10,000 iterations).

\subsection{SciERC Ablation}
\label{sec:ablation}

\begin{table}[t]
\centering
\small
\begin{tabular}{lcccccc}
\toprule
\textbf{Config} & \textbf{DR\%} & \textbf{ET\%} & \textbf{RV\%} & \textbf{CF\%} & \textbf{E-F1} & \textbf{R-F1} \\
\midrule
Unconstr. & 15.37 & 0.40 & 1.39 & 0.08 & 31.01 & 8.23 \\
Gram. & 14.84 & 0.00 & 0.00 & 0.08 & 30.98 & 8.23 \\
Sem. & 11.65 & 0.40 & 1.51 & 2.71 & 31.01 & 8.15 \\
Full S1 & \textbf{0.93} & 0.00 & 0.00 & 9.90 & 30.98 & \textbf{8.36} \\
Full S2 & 0.97 & 0.00 & 0.00 & 8.03 & 31.90 & 8.03 \\
\bottomrule
\end{tabular}
\caption{SciERC ablation (Qwen3-8B). Full S1 achieves 94\% DR reduction. Grammar eliminates ET and RV; semantic adds DR enforcement. Semantic-only achieves only 24\% DR reduction---the enabling dependency in action.}
\label{tab:scierc-ablation}
\end{table}

Table~\ref{tab:scierc-ablation} presents the SciERC ablation.
The full S1 configuration reduces DR from 15.37\% to 0.93\% ($\downarrow$94\%) while maintaining F1-neutral impact: R-F1 $\Delta{=}{+}0.13$pp ($p{=}0.58$, 95\% CI [$-$0.32, $+$0.58]).

The ablation cleanly separates the two layers' roles:
\begin{itemize}
\item \textbf{Grammar-only}: ET and RV $\to$ 0\%, but DR barely changes (14.84\%, $\Delta{=}{-}0.53$pp). Grammar ensures valid vocabularies but does not understand type semantics.
\item \textbf{Semantic-only}: DR drops to 11.65\% ($\downarrow$24\%), far below Full S1's 94\%. With $\sim$1.4\% of relations out-of-vocabulary, the semantic layer cannot look up constraints for those triples.
\item \textbf{Full S1}: The combination achieves $\downarrow$94\%---grammar enables semantic by guaranteeing valid relation names for constraint lookup.
\end{itemize}

The residual DR of 0.93\% (16 triples out of 1,729) corresponds to category~(c) constraints in our decidability taxonomy: truly non-local constraints requiring forward coreference resolution, which is beyond single-pass forward-only decoding.

\paragraph{S2 as Architectural Ablation.}
S2 achieves comparable DR (0.97\%) but slightly lower R-F1 (8.03\% vs.\ 8.36\%) and does not meaningfully reduce CF (8.03\% vs.\ 9.90\%, $\Delta{=}{-}1.87$pp).
We retain S1 as the recommended configuration.

\subsection{DocRED Results}

\begin{table}[t]
\centering
\small
\begin{tabular}{lcccccc}
\toprule
\textbf{Config} & \textbf{DR\%} & \textbf{ET\%} & \textbf{RV\%} & \textbf{CF\%} & \textbf{E-F1} & \textbf{R-F1} \\
\midrule
Unconstr. & 18.46 & 2.48 & 29.66 & 0.10 & 58.15 & 7.74 \\
Gram. & 26.17 & 0.00 & 0.03 & 0.07 & 59.38 & 8.31 \\
Sem. & 13.84 & 4.22 & 29.87 & 1.05 & 57.81 & 7.70 \\
Full S1 & \textbf{10.21} & 0.00 & 0.03 & 2.57 & \textbf{59.37} & \textbf{8.23} \\
\bottomrule
\end{tabular}
\caption{DocRED ablation (Qwen3-8B). Full S1 achieves 44.7\% DR reduction with significant F1 improvement (R-F1 $+$0.50pp, $p{<}0.001$; E-F1 $+$1.23pp, $p{<}0.001$). Grammar-only \emph{increases} DR---see text for explanation.}
\label{tab:docred-ablation}
\end{table}

Table~\ref{tab:docred-ablation} shows DocRED results.
Full S1 reduces DR from 18.46\% to 10.21\% ($\downarrow$44.7\%) with \emph{significant F1 improvement}: R-F1 $+$0.50pp ($p{<}0.001$, 95\% CI [$+$0.23, $+$0.81]), E-F1 $+$1.23pp ($p{<}0.001$, 95\% CI [$+$0.94, $+$1.55]).
The F1 gains are attributable entirely to grammar-layer OOV elimination: grammar-only E-F1 (59.38\%) $\approx$ Full S1 E-F1 (59.37\%), confirming the semantic layer contributes DR reduction but not F1 improvement.

\paragraph{Grammar-Only DR Increase.}
A counterintuitive finding: grammar-only \emph{increases} DR from 18.46\% to 26.17\%.
When grammar forces OOV relation names (29.66\% of triples) to map to valid vocabulary entries, the randomly assigned valid names often have domain/range constraints incompatible with the generated entity types.
Grammar \emph{exposes latent violations} that were previously hidden behind OOV names (which lack defined constraints).
This is the key insight for understanding the interaction between vocabulary and semantic constraint layers.

\paragraph{Tightness $\times$ RV Baseline Interaction.}
The reduced DR elimination on DocRED (45\% vs.\ 94\% on SciERC) is explained by a dual-factor model: Qwen's high RV baseline on DocRED (29.66\%) causes grammar to expose many latent violations, inflating the effective DR that the semantic layer must address.
Cross-model validation confirms this: Llama-3.1-8B, with a lower RV baseline (16.03\%) on DocRED, achieves 90.2\% DR reduction on the same ontology (\S\ref{sec:crossmodel}).

\subsection{Cross-Model Validation}
\label{sec:crossmodel}

\begin{table}[t]
\centering
\small
\begin{tabular}{llcccc}
\toprule
\textbf{Model} & \textbf{Dataset} & \textbf{DR\textsubscript{unc}} & \textbf{DR\textsubscript{S1}} & \textbf{$\downarrow$\%} & \textbf{$\Delta$R-F1} \\
\midrule
\multirow{2}{*}{Qwen} & SciERC & 15.37 & 0.93 & 94.0 & $+$0.13 \\
& DocRED & 18.46 & 10.21 & 44.7 & $+$0.50$^{***}$ \\
\midrule
\multirow{2}{*}{Llama} & SciERC & 13.34 & 0.08 & 99.4 & $-$0.14 \\
& DocRED & 26.07 & 2.55 & 90.2 & $+$0.29$^{*}$ \\
\bottomrule
\end{tabular}
\caption{Cross-model $\times$ cross-dataset DR reduction summary. All experiments use zero code modification across models. $^{***}$: $p{<}0.001$; $^{*}$: $p{<}0.05$.}
\label{tab:crossmodel}
\end{table}

Table~\ref{tab:crossmodel} presents the full 2-model $\times$ 2-dataset matrix.
\sys{} achieves substantial DR reductions (45--99\%) across all configurations with zero code modification.

Llama-3.1-8B on SciERC achieves 99.4\% DR reduction (13.34\%$\to$0.08\%, 1 residual violation out of 1,315 triples), approaching the theoretical limit of forward-only decoding.
On DocRED, Llama achieves 90.2\% reduction (26.07\%$\to$2.55\%).

The variation in reduction rates is fully explained by the tightness $\times$ RV baseline interaction:
Qwen's high RV on DocRED (29.66\%) causes grammar to expose many latent DR violations, yielding a lower relative reduction (45\%).
Llama's lower RV (16.03\%) means fewer latent violations are exposed, enabling 90.2\% reduction.
The method itself works consistently; observed variation reflects model-specific vocabulary compliance, not constraint enforcement failure.

An independent finding emerges from the cross-model comparison: Llama exhibits \emph{orthogonal} RV and DR profiles compared to Qwen on DocRED.
Llama has lower RV (16.03\% vs.\ 29.66\%) but higher DR (26.07\% vs.\ 18.46\%).
This demonstrates that vocabulary compliance and semantic compliance are genuinely independent failure modes, empirically justifying the two-layer architecture.

\subsection{Comparison with Post-Hoc Filtering}
\label{sec:posthoc}

\begin{table}[t]
\centering
\small
\begin{tabular}{lcccc}
\toprule
& \multicolumn{2}{c}{\textbf{SciERC}} & \multicolumn{2}{c}{\textbf{DocRED}} \\
\cmidrule(lr){2-3} \cmidrule(lr){4-5}
& \sys{} & Post-hoc & \sys{} & Post-hoc \\
\midrule
R-F1 & 8.36 & 8.63 & 8.23 & 9.61 \\
Recall $\Delta$ & $+$0.20 & $-$0.72 & $+$0.56 & $-$1.06 \\
Reject. rate & 0\% & 16.8\% & 0\% & 48.1\% \\
TP lost & 0 & --- & 0 & 127 \\
ET residual & 0\% & 0.40\% & 0\% & 2.47\% \\
\bottomrule
\end{tabular}
\caption{Multi-dimensional comparison with post-hoc filtering. Post-hoc achieves higher precision by discarding 17--48\% of triples. \sys{} preserves all triples, improves recall, and fully eliminates entity type violations.}
\label{tab:posthoc}
\end{table}

Table~\ref{tab:posthoc} compares \sys{} with post-hoc filtering (removing all violated triples from unconstrained output).
Post-hoc achieves higher R-F1 by aggressively removing false positives, but at substantial cost:
(1)~\textbf{Recall loss}: post-hoc recall drops ($-$0.72pp SciERC, $-$1.06pp DocRED) while \sys{} recall increases ($+$0.20pp, $+$0.56pp) via probability mass redistribution to valid tokens.
(2)~\textbf{Information loss}: post-hoc discards 48.1\% of triples on DocRED---nearly half the extracted information. On DocRED, 127 true-positive triples are lost because they were correctly extracted but had entity type violations that \sys{} prevents at generation time.
(3)~\textbf{Incomplete enforcement}: post-hoc cannot address entity type violations (residual ET 2.47\% on DocRED), as it only filters triples, not entity lists. \sys{} achieves ET$=$0\%.
(4)~\textbf{Operational}: \sys{} produces constraint-compliant output in a single pass suitable for streaming; post-hoc requires an additional verification pipeline.

\subsection{Statistical Rigor}

\begin{table}[t]
\centering
\small
\begin{tabular}{llccc}
\toprule
\textbf{Setting} & \textbf{Metric} & \textbf{$\Delta$} & \textbf{95\% CI} & \textbf{$p$} \\
\midrule
Qwen SciERC & R-F1 & $+$0.13 & [$-$0.32, $+$0.58] & 0.58 \\
Qwen DocRED & R-F1 & $+$0.50 & [$+$0.23, $+$0.81] & ${<}$0.001 \\
Qwen DocRED & E-F1 & $+$1.23 & [$+$0.94, $+$1.55] & ${<}$0.001 \\
Llama SciERC & R-F1 & $-$0.14 & --- & 0.011 \\
Llama DocRED & R-F1 & $+$0.29 & --- & 0.013 \\
\bottomrule
\end{tabular}
\caption{Bootstrap confidence intervals (10,000 iterations) for all F1 claims. SciERC R-F1 impact is F1-neutral ($p{=}0.58$). DocRED R-F1 and E-F1 improvements are significant. Llama SciERC R-F1 decline ($-$0.14pp) is statistically detectable but practically negligible.}
\label{tab:bootstrap}
\end{table}

Table~\ref{tab:bootstrap} reports bootstrap confidence intervals for all F1 claims.
We distinguish statistical significance from practical relevance: the Llama SciERC R-F1 decline ($-$0.14pp, $p{=}0.011$) is statistically detectable but represents a 5\% relative change on a 2.68\% baseline---practically negligible.

\subsection{Cross-Field Consistency Analysis}
\label{sec:cf-analysis}

\sys{} introduces elevated cross-field consistency violations (CF): 9.90\% on SciERC Full S1 (vs.\ 0.08\% unconstrained) and 2.57--12.02\% on DocRED depending on model.
Ablation confirms CF is introduced entirely by the semantic layer: grammar-only CF $=$ 0.08\% (unchanged from baseline).

The mechanism is \emph{entity-vs-relation type drift} (EVR): the semantic layer enforces domain/range constraints independently per relation, which can assign different types to the same entity across different relation contexts.
Analysis of the dominant conflict patterns on SciERC reveals that $>$60\% involve \textsc{Generic}$\leftrightarrow$\textsc{OtherScientificTerm} confusion---type pairs with genuinely ambiguous boundaries in the SciERC type system.

CF violations represent a qualitatively different failure mode than DR violations: each individual triple remains type-valid (domain/range satisfied), only cross-triple consistency is broken.
DR violations produce structurally invalid triples unusable by downstream systems; CF violations produce individually valid triples with inconsistent entity typing across contexts.

\section{Related Work}
\label{sec:related}

\paragraph{Constrained Decoding.}
Grammar-constrained decoding enforces structural constraints during generation \cite{Willard2023Outlines,Dong2024XGrammar}.
SynCode \cite{Ugare2024SynCode} extends grammar constraints to incremental parsing.
DOMINO \cite{Beurer2024DOMINO} combines grammar and semantic constraints for data modeling.
\sys{} differs in targeting IE-specific ontology constraints with a two-layer architecture exhibiting principled enabling dependency.

\paragraph{LLM-based Information Extraction.}
Recent work demonstrates strong zero-shot IE capabilities of LLMs \cite{Wei2023ZeroShotIE,Wadhwa2023Revisiting,Li2023EvaluatingChatGPT} but does not systematically address ontology constraint compliance.
GoLLIE \cite{Sainz2024GoLLIE} uses code-style annotation guidelines to improve extraction quality.
\sys{} is complementary: it enforces ontology constraints regardless of prompt design or model capability.

\paragraph{Type Constraint Enforcement.}
Knowledge graph embedding methods incorporate type constraints as soft regularizers \cite{Krompass2015TypeKG}.
\sys{} provides hard constraint enforcement during generation, guaranteeing compliance rather than encouraging it.

\section{Conclusion}

We presented a systematic study of ontology constraint violations in LLM-based IE, finding that 13--26\% of triples violate domain/range constraints and that violation rates are driven by constraint tightness rather than ontology size.
Schema-informed prompting is fundamentally unreliable---it can increase violations for specific relation types.

\sys{} compiles ontology schemas into composable decoding-time constraints, achieving 45--99\% domain/range violation reduction across two datasets and two model families at zero or positive F1 cost, with zero code modification across models.
The key architectural insight is the enabling dependency between grammar and semantic layers: grammar enforcement is a prerequisite for effective semantic constraint enforcement, not merely an additive component.

\paragraph{Limitations.}
\sys{} introduces cross-field consistency violations (CF up to 12\%) due to per-relation independent type enforcement.
This is an inherent limitation of forward-only decoding without global entity-type consensus.
Future work should explore a cross-relation consistency pass to resolve entity-level type conflicts.
Additionally, our evaluation is limited to zero-shot settings with two model families; validation on fine-tuned models and additional ontologies would strengthen generalization claims.

\bibliographystyle{plainnat}
\begin{thebibliography}{20}

\bibitem[Beurer-Kellner et~al.(2024)]{Beurer2024DOMINO}
Beurer-Kellner, L., Fischer, M., and Vechev, M.
\newblock DOMINO: Domain-aware model calibration in medical pathology.
\newblock In \emph{ICML}, 2024.

\bibitem[Dong et~al.(2024)]{Dong2024XGrammar}
Dong, Y., Huang, C., and Chen, T.
\newblock XGrammar: Flexible and efficient structured generation engine for large language models.
\newblock \emph{arXiv preprint arXiv:2411.15100}, 2024.

\bibitem[Krompass et~al.(2015)]{Krompass2015TypeKG}
Krompass, D., Baier, S., and Tresp, V.
\newblock Type-constrained representation learning in knowledge graphs.
\newblock In \emph{ISWC}, 2015.

\bibitem[Li et~al.(2023)]{Li2023EvaluatingChatGPT}
Li, B., Fang, G., Yang, Y., Wang, Q., Ye, W., Zhao, S., and Zhang, R.
\newblock Evaluating ChatGPT's information extraction capabilities.
\newblock \emph{arXiv preprint arXiv:2304.11633}, 2023.

\bibitem[Luan et~al.(2018)]{Luan2018SciERC}
Luan, Y., He, L., Ostendorf, M., and Hajishirzi, H.
\newblock Multi-task identification of entities, relations, and coreference for scientific knowledge graph construction.
\newblock In \emph{EMNLP}, 2018.

\bibitem[Sainz et~al.(2024)]{Sainz2024GoLLIE}
Sainz, O., de~Lacalle, O.~L., Labaka, G., Barrena, A., and Agirre, E.
\newblock GoLLIE: Annotation guidelines improve zero-shot information extraction.
\newblock In \emph{ICLR}, 2024.

\bibitem[Ugare et~al.(2024)]{Ugare2024SynCode}
Ugare, S., Suresh, T., Kang, H., Misailovic, S., and Singh, G.
\newblock SynCode: LLM generation with grammar augmentation.
\newblock In \emph{NeurIPS}, 2024.

\bibitem[Wadhwa et~al.(2023)]{Wadhwa2023Revisiting}
Wadhwa, S., Amir, S., and Wallace, B.~C.
\newblock Revisiting relation extraction in the era of large language models.
\newblock In \emph{ACL}, 2023.

\bibitem[Wei et~al.(2023)]{Wei2023ZeroShotIE}
Wei, X., Cui, X., Cheng, N., Wang, X., Zhang, X., Huang, S., Xie, P., Xu, J., Chen, Y., Zhang, M., Jiang, Y., and Han, W.
\newblock Zero-shot information extraction via chatting with ChatGPT.
\newblock \emph{arXiv preprint arXiv:2302.10205}, 2023.

\bibitem[Willard and Louf(2023)]{Willard2023Outlines}
Willard, B.~T. and Louf, R.
\newblock Efficient guided generation for large language models.
\newblock \emph{arXiv preprint arXiv:2307.09702}, 2023.

\bibitem[Yao et~al.(2019)]{Yao2019DocRED}
Yao, Y., Ye, D., Li, P., Han, X., Lin, Y., Liu, Z., Liu, Z., Huang, L., Zhou, J., and Sun, M.
\newblock DocRED: A large-scale document-level relation extraction dataset.
\newblock In \emph{ACL}, 2019.

\end{thebibliography}

\end{document}
